\documentclass[12pt, a4paper]{article}

\usepackage{fontspec}
\usepackage{geometry}
\usepackage{amsmath, amssymb}
\usepackage{microtype}
\usepackage{setspace}
\usepackage{titlesec}
\usepackage{hyperref}
\usepackage{xcolor}
\usepackage{booktabs}
\usepackage{amsthm}

\newtheorem{theorem}{Theorem}
\newtheorem{corollary}{Corollary}[theorem]

\geometry{margin=1.25in}
\setstretch{1.35}

\definecolor{darkgray}{HTML}{2C2C2C}
\color{darkgray}

\titleformat{\section}{\large\bfseries}{{\thesection.}}{0.6em}{}
\titleformat{\subsection}{\normalsize\bfseries\itshape}{}{0em}{}
\titlespacing*{\section}{0pt}{2.5ex plus 0.5ex minus 0.2ex}{1ex plus 0.2ex}
\titlespacing*{\subsection}{0pt}{1.5ex plus 0.3ex}{0.5ex}

\hypersetup{
  colorlinks=true,
  linkcolor=black,
  urlcolor=black,
}

\setlength{\parindent}{1.5em}
\setlength{\parskip}{0.4ex}

\title{\textbf{The Coordinate-System Error:\\
Misdiagnosis in the Theories of Attention, Efficiency, and Technological Progress}}
\author{Flyxion}
\date{\today}

\begin{document}

\maketitle

\begin{abstract}
\noindent
Two registers of contemporary discourse share a common structural defect. The first, exemplified by accounts of continuous partial attention such as that advanced by Tom Scryleus, diagnoses modern cognitive life as a condition of perpetual fragmentation. The second, exemplified by industry advocates of artificial intelligence, diagnoses the same technological transformation as a condition of expanding efficiency and democratized access. Both diagnoses are rendered in a coordinate system that is inadequate to the phenomena it purports to describe. In each case, a scalar, single-dimensional framework is applied to processes that are irreducibly multi-layered, field-like, and trajectory-dependent. The result is a consistent pattern of misinterpretation: behaviors that are coherent within a higher-dimensional representation appear as pathology or as progress in the impoverished one, depending only on whether the projection discards the costs or the capabilities.

This essay argues that the correct object of analysis is neither the instantaneous state of attention nor the internal efficiency metrics of a technological system, but the trajectory of irreversible events within a persistent, verifiable historical field. Drawing on personal practice across construction, language acquisition, and sustained multi-source learning, and on a structural analysis of the economic and epistemic consequences of current AI deployment, the essay shows that both the fragmentation thesis and the efficiency thesis commit the same error at different scales: they suppress irreversibility, externalize cost, and mistake the limits of their own representational medium for facts about the world.

The argument is developed in three movements. The first overturns the scalar attention model and replaces it with a field-stratified account of cognition and parallel constraint resolution. The second extends the same diagnosis into AI and economic discourse, showing that efficiency metrics are formally non-injective projections that necessarily accumulate informational debt. The third grounds both movements in a unified framework of irreversibility, entropy, and the plenum, formalized through a set of explicit primitives and a central theorem from which both domains follow as corollaries. Additional sections address the isomorphism between cognitive and infrastructural misdiagnosis, the geometry of informational chokepoints, the scalar assessment trap in pedagogy, the forensic fragility produced by synthetic media, the structural persistence of the error under economic conditions, and the design criteria that a genuinely adequate system must satisfy. The solution in every case is not a change of policy but an expansion of the dimensionality of the models used to describe the system \cite{Prigogine1984,Shannon1948,Friston2010}.
\end{abstract}

\vspace{1em}

\section{Introduction: A Misdiagnosis at the Level of the Model}

Scryleus begins from a phenomenological observation that is not in dispute: something has changed in how people occupy shared physical space. At bus stops, dinner tables, and elevators, bodies are co-located while cognition is distributed across multiple relational contexts simultaneously. He names this condition ``continuous partial attention,'' distinguishes it from ordinary multitasking, and argues that it is not a strategy for efficiency but a structural condition of modern life in which the self is never fully anywhere.

The observations are real. The interpretation is systematically mistaken, and the error is not incidental. It arises from a foundational premise that Scryleus never makes explicit: that presence is a scalar quantity localized in physical space and that attention is a single-threaded, bandwidth-limited resource that diminishes whenever it is divided \cite{Simon1969}. Once that premise is in place, every observation confirms the diagnosis. A person on a phone is ``absent.'' Multiple simultaneous inputs constitute ``fragmentation.'' Notifications ``move attention before choice.'' Social contracts dissolve because no one is ``fully present enough to hold them.'' The argument is internally coherent, but it is coherent in the way a map is coherent when the projection is wrong: every feature is distorted in the same direction, and the distortion is invisible from within the map.

The scalar model and the field-stratified model presented in this essay are not merely competing descriptions of the same phenomena. They are non-equivalent representations that preserve different invariants. Under the scalar model, the conserved quantity is the amount of attention allocated to a single locus. Under the field-stratified model, the conserved quantity is coherence across trajectories over time. Because the models preserve different quantities, the same observations produce opposite conclusions---and because certain dynamics of layered cognition have no representation in the scalar system, behaviors that are perfectly regular in the richer framework must appear as anomalies or pathologies in the impoverished one \cite{Hutchins1995,Clark2013}.

This essay develops the argument across two movements. The first addresses the continuous partial attention thesis directly, replacing each of its impoverished premises with a structurally adequate one. The second extends the same diagnosis to the dominant discourse surrounding artificial intelligence, showing that the efficiency thesis commits an identical error at the scale of economic and infrastructural systems. Both movements converge on a single conclusion: the problem is not fragmentation, not inefficiency, and not lack of technological access. The problem is that the models currently used to evaluate these systems lack the degrees of freedom necessary to represent what is actually taking place.

\section{Presence Is a Field, Not a Point}

\subsection{The claim}
Scryleus's foundational charge is that contemporary individuals are ``physically here, mentally elsewhere.'' Presence has shifted into a condition of perpetual absence.

\subsection{The counter}
This characterization assumes that presence is singular and spatially localized---that co-location and full presence are, or ought to be, the same thing. Because co-location and attention once reliably coincided, the loss of that coincidence is being read as the loss of presence itself.

The more precise description is that presence has changed its topology. A person simultaneously engaged in a face-to-face conversation, a text exchange with a geographically distant colleague, and background monitoring of a developing situation is not absent from any of these contexts: they are multiply embedded across several relational manifolds at once. Attention is not a point; it is closer to a vector field whose support can span many contexts without being fully absent from any of them \cite{Friston2010}. The observation that single-location presence no longer exhausts the structure of engagement is accurate. The inference that distributed presence is equivalent to no presence does not follow.

This topological shift is genuinely novel and requires new frameworks to describe. Scryleus does not develop such frameworks. He applies the older scalar model, observes that presence no longer collapses to a single spatial coordinate, and concludes that presence has been lost. The gap between observation and conclusion is precisely where the coordinate-system error is introduced.

\section{Parallel Constraint Resolution Versus Attentional Bandwidth}

\subsection{The claim}
Continuous partial attention is distinguished from productive multitasking on the grounds that it fragments cognition without generating the efficiencies that genuine multitasking promises. The divided mind produces shallow outputs from each stream.

\subsection{The counter}
This distinction rests on a hidden model: that cognition is a serial pipeline with a bounded attention budget, such that allocation to one task necessarily depletes allocation to another \cite{Simon1969}. Under this model, switching between streams incurs a context-reconstruction cost, and the aggregate result is superficial engagement with each.

This model is correct under a specific and restricted regime: when tasks are independent and compete for the same representational substrate. It fails when tasks are correlated projections of a shared latent structure. In the latter case, moving between sources is not context switching---it is constraint propagation. Inconsistencies between representations generate error signals that guide attention toward unresolved structure. Redundancies confirm and reinforce what has already been acquired. Switching accelerates convergence rather than disrupting it \cite{Clark2013}.

The formal contrast is between
\[
\text{Task}_A + \text{Task}_B \;\rightarrow\; \text{interference},
\]
which holds for independent, substrate-competing streams, and
\[
\text{Motor field} \;\oplus\; \text{Semantic field} \;\rightarrow\; \text{compositional throughput},
\]
which holds when streams occupy heterogeneous representational layers with structural relations between them.

This is not an abstract theoretical claim but a constructive counterexample to the universal thesis that divided attention reduces performance and depth. My own practice during several years of construction work makes this concrete. Demolition, wiring, and framing engage sensorimotor systems that, once trained, operate with procedural autonomy: they do not continuously demand the representational resources that semantic comprehension requires. Running university lectures through headphones simultaneously is not splitting a single attentional budget; it is operating on multiple layers in parallel, with the motor layer exhibiting the procedural independence that permits the semantic layer to process without interference \cite{Hutchins1995,Suchman2007}. The result was the compression of approximately seven years of lecture content into eighteen months, while maintaining execution quality across physically demanding skilled-trade work. That is not fragmentation. It is integration across strata.

Because this is a constructive counterexample to a universal claim---the claim that divided attention inherently reduces depth---the burden of proof shifts completely. The continuous partial attention thesis cannot be universally true. It must be restricted to the narrower regime in which streams are genuinely independent and substrate-competing, and that restriction is precisely what the structural analysis provides.

\section{Agency in a Modified Potential Field}

\subsection{The claim}
Attention is engineered and bypassed: notifications, infinite scrolls, and autoplay move attention before a choice is made, constituting a form of optimization that operates below the threshold of conscious deliberation and effectively removes agency from the individual.

\subsection{The counter}
The observation that these systems create high-gradient attractors is correct. The inference that this eliminates agency is not. The distinction between loss of agency and navigation within a modified potential field is not merely rhetorical: loss of agency implies deterministic capture, whereas a modified field implies biased navigation with preserved degrees of freedom. Only the second is consistent with observed behavior.

In the modified-field picture, attention follows gradient descent in an artificially restructured energy landscape \cite{Friston2010}. The system does not remove agency---it raises the cost of certain choices while reducing the cost of others. Reaching for a phone during a lull carries near-zero cost in the current landscape. Maintaining attention on an unstructured internal state against a high-gradient attractor carries significant cost. Agency is not absent; it is being exercised under asymmetric cost conditions that have been deliberately engineered for commercial purposes.

This distinction matters because it changes what the solution looks like. If agency is absent, the prescription must be dispositional: cultivate stronger will, resist impulse, reclaim intentionality. If the problem is structural, the prescription is redesign: modify the landscape so that gradient directions align with intended trajectories rather than against them \cite{Zuboff2019}. Behavioral adaptation that restructures the environment---removing devices from certain spaces, modifying notification architecture, designing deliberate friction into high-gradient attractors---is already implementing this solution at a small scale. The correct framing is not that one has ceased to choose, but that the cost of certain choices has been artificially increased by systems designed to exploit attentional dynamics for commercial ends.

\section{Phenomenological Thinness as a Loss of Representational Degrees of Freedom}

\subsection{The claim}
Constant stimulation produces a ``thinness'' of experience: time feels shallow, engagement lacks depth, and the qualitative richness of immediate experience is sacrificed to a baseline of perpetual low-level stimulation.

\subsection{The counter}
The phenomenological observation is genuine. The explanation is incorrect, and the correction requires precision about where the loss occurs.

Two-dimensional, sequential interfaces impose a specific and severe compression on cognitive processes that are inherently high-dimensional. They linearize experience, erase spatial persistence, and collapse relational structure into flat sequences. This compression produces a subjective loss of depth---not because cognition has become shallower, but because the representational medium has lost the degrees of freedom required to preserve the depth that is present. What is being experienced is not a reduction in cognitive capacity but a loss of representational degrees of freedom under dimensional compression \cite{Landauer1961,Jaynes1957}. The thinness is an artifact of the mapping from cognitive space to interface space, not a property of the cognitive space itself.

The evidence that this is a representational rather than cognitive fact is behavioral. My earlier academic practice---reading multiple textbooks simultaneously in physical space, working on several courses at once while taking notes, drawing diagrams across spread-out surfaces, moving between sources by spatial proximity---was producing the same high-throughput, multi-stream processing that Scryleus identifies as the problem. The difference was that the physical medium allowed spatial indexing, persistent partial states, and relational layout. Books could be positioned by associative proximity. Diagrams could maintain intermediate structure across a session. The environment encoded part of the cognitive work, functioning as an external representational scaffold that preserved dimensionality \cite{Hutchins1995}.

Digital interfaces strip most of this away. The result looks like attentional degradation because the scaffolding is absent, and the degree-of-freedom budget of the interface is insufficient to represent the relational structure that physical arrangement previously maintained. Given the right substrate---physical, spatial, persistent, semantically indexed---the same patterns of engagement that read as distraction on a screen reveal what they actually are: efficient navigation of a high-dimensional cognitive field. The mismatch is between cognitive geometry and interface geometry, not between current cognition and some prior, superior form of it.

\section{Temporal Evaluation and the Illusion of Fragmentation}

\subsection{The claim}
Continuous partial attention appears as fragmentation because, at any given instant, attention is visibly distributed across multiple streams without full resolution in any single one of them.

\subsection{The counter}
The appearance of fragmentation is an artifact of evaluating cognition at the wrong temporal scale. The analysis underlying the continuous partial attention thesis operates on instantaneous slices, effectively sampling the system at time $t$ and inferring structure from the state alone. Under such sampling, any system engaged in parallel processing will appear incoherent, because coherence in a distributed system is not a property of the instantaneous state but of the trajectory \cite{Wiener1948,Ashby1956}.

The correct object of evaluation is not the state at time $t$ but the trajectory over the interval $[0, T]$. At the level of trajectory, patterns that appear disjoint in the moment resolve into coherent structure. Language acquisition through distributed flashcard practice, lecture intake interleaved with physical labor, and multi-source reading that builds cross-referenced understanding all exhibit this property: no single moment contains the full structure, but the accumulated history does. Coherence is longitudinal, not instantaneous.

The illusion of fragmentation arises directly from this mismatch between the temporal resolution of observation and the temporal scale at which coherence emerges. This is not a minor technical point. It is the reason that constructive learning practices---thirty thousand Arabic flashcard items, compressed lecture intake, the university practice of reading across multiple books simultaneously---could look from the outside like exactly the behavior being condemned, while producing precisely the outcomes the critique claims that behavior prevents. The behavior is identical at the level of the instantaneous sample. The dynamics are entirely different at the level of the trajectory.

\section{Social Enforcement Under Cognitive Load}

\subsection{The claim}
Social norms erode because people are no longer ``present enough'' to hold shared agreements. The bus-stop example---where the unspoken rules of queuing and eye contact go unenforced---stands for a broader dissolution of the social fabric.

\subsection{The counter}
Social norm enforcement is a resource-dependent process. It requires attention, situational evaluation, and willingness to intervene. In any environment, enforcement is probabilistic and subject to allocation. The critical distinction here is between capacity and allocation: nothing in the observation implies reduced capacity to enforce norms. What has changed is the priority weighting of enforcement relative to other cognitive demands.

In an environment of increased total cognitive load, local enforcement probability drops as a predictable consequence of demand exceeding the available surplus. Formally,
\[
P(\text{enforcement}) \;\propto\; \text{available attentional surplus}.
\]
Increasing global load decreases local enforcement probability without implicating any change in the underlying cognitive capacities involved in norm recognition, evaluation, or intervention \cite{Ashby1956}. People retain the full capacity to notice and respond to norm violations. They allocate that capacity differently under increased total demand. What Scryleus reads as absence is more precisely reweighted cost functions operating under changed conditions. The social fabric is not dissolving; enforcement of low-priority norms is being selectively suppressed in a high-load environment, exactly as one would expect from any resource-constrained allocation system.

\section{Stress, Urgency, and the Expansion of Option Space}

\subsection{The claim}
Digital fragmentation produces increased stress and a persistent low-level sense of urgency that pervades everyday experience.

\subsection{The counter}
The phenomenological correlation is real. The causal attribution is imprecise. Stress of the kind described is better modeled as a function of the size of the available option space. Letting $\Omega_t$ denote the set of available actions, accessible contexts, and unresolved commitments at time $t$,
\[
\text{Stress} \;\sim\; |\Omega_t|.
\]
Mobile devices and persistent connectivity have dramatically expanded $|\Omega_t|$. Every moment now contains more possible actions, more accessible informational contexts, and more open commitments whose status can be checked or revised. This produces persistent background obligation---not because attention has been damaged, but because the decision surface has grown and the option space cannot be fully collapsed \cite{Simon1969}.

The key additional observation is that the same structural change that produces stress also produces unprecedented latent agency. An expanded $|\Omega_t|$ is not simply a burden; it is also the condition of possibility for the epistemic access, cross-domain synthesis, and distributed collaboration that were previously unavailable to most individuals. The stress and the capability are two faces of the same structural transformation. Treating the stress in isolation, as the critique does, misrepresents both the phenomenon and the appropriate response to it.

\section{Polarization as Phase Separation}

\subsection{The claim}
Digital fragmentation deepens social polarization, producing reductive thinking because distracted individuals can no longer hold the complexity required for nuanced engagement.

\subsection{The counter}
Polarization does not require reduced cognitive capacity, and the causal argument for attentional fragmentation as its source is one of the weakest components of the thesis. Higher connectivity enables finer clustering, faster within-group alignment, and cleaner between-group boundary formation. What appears as division is frequently phase separation under increased informational resolution---the system revealing latent structure that was always present but previously suppressed by the friction and blur of lower-bandwidth information exchange \cite{Prigogine1984}.

Moreover, polarization has many independent determinants---economic incentives structuring media production, affective asymmetries in social reinforcement dynamics, the political economy of attention itself---whose contributions are larger and more direct than those of attentional style \cite{Zuboff2019}. Attributing the phenomenon primarily to fragmented attention conflates a real but limited effect with a structural problem whose deeper causes lie elsewhere.

\section{Creation, Consumption, and Trajectory Formation}

\subsection{The claim}
The act of making or creating serves as the necessary antidote to fragmented attention, because creative work demands a kind of presence that passive consumption destroys.

\subsection{The counter}
This is the most structurally accurate observation in the text, but it is misinterpreted as a moral rather than a formal distinction. The relevant difference is not between virtuous creation and degrading consumption but between two attentional regimes with different structural properties.

Trajectory formation involves irreversible commitment to a developing course of action \cite{Landauer1961}. Each decision constrains the next; coherence must be maintained across the whole; the option space progressively collapses as the work takes form. Consumption without accumulative structure involves movement through possible states without convergence: the option space remains open or expands, and nothing collapses. Creation forces the first regime because output is irreversible. This is why it demands the particular quality of attention Scryleus identifies---not because it is morally superior, but because the task structure requires sustained coherence along a trajectory that cannot be undone.

Both regimes are necessary, and neither is intrinsically degrading. The problem is not consumption per se but consumption structured so as to increase entropy without resolution: streams that expand the option space without ever closing it \cite{Jaynes1957}. The appropriate intervention is not to moralize against consumption but to design consumption environments that maintain accumulative structure---to build the system so that attending to inputs contributes to the formation of a long-term trajectory rather than diffusing into an undifferentiated background.

\section{Embodiment and the Partition of Cognitive Load}

\subsection{The claim}
Simultaneous engagement with multiple streams necessarily divides attention and reduces performance.

\subsection{The counter}
This claim presupposes that all cognitive tasks draw from a single undifferentiated resource pool. In practice, cognitive architecture is partitioned across heterogeneous subsystems with differing bandwidth constraints and degrees of functional autonomy. Sensorimotor processes, once trained, operate with a high degree of procedural independence from the top-down control processes that govern semantic comprehension. Semantic processing in the auditory domain can proceed in parallel with skilled physical execution because the two engage distinct representational substrates and do not compete for the same layer of cognitive resources \cite{Clark2013,Suchman2007}.

The relevant distinction is therefore not between single-task and multi-task engagement but between tasks that compete for the same representational substrate and tasks that occupy distinct but interacting layers of the cognitive hierarchy. In the former case, interference is real and performance degrades. In the latter, parallel engagement increases effective throughput without degradation in either layer. The combination of skilled physical labor with sustained lecture intake exemplifies the second case and is not an exceptional individual achievement but a consequence of the layered architecture of embodied cognition under appropriate training conditions \cite{Hutchins1995}.

\section{The Infinite Scroll Inversion}

There is an autobiographical dimension to this argument with direct formal relevance. Two decades before smartphones, I produced approximately twenty thousand Spanish flashcards, practised them in fragments of available time, and accumulated vocabulary over months of distributed rehearsal. Later, I completed approximately thirty thousand items on Memrise to acquire reading and writing fluency in Arabic. Viewed from the outside, both practices are structurally isomorphic to the behavior Scryleus identifies as the pathological form: small units of engagement, repeated brief exposure, interleaving across multiple contexts, utilization of otherwise idle time, accumulation oriented toward a long horizon.

The interaction pattern is identical. The cognitive outcome is entirely different. The decisive variable is whether the stream is entropy-reducing or entropy-increasing \cite{Shannon1948,Jaynes1957}. In a structured accumulative regime, fragments are organized, integration is tracked, each element reduces uncertainty, and there exists a convergence criterion against which progress can be measured. In the unstructured regime typical of optimized social feeds, fragments are selected for affective gradient rather than coherence, accumulation is not integrated into any long-term structure, entropy is injected without resolution, and no convergence criterion exists.

The formal equivalence runs precisely the other way from the critique's assumption: the behavior being condemned and the behavior that produced deep language acquisition and procedural fluency are isomorphic at the level of interaction pattern, differing only in whether the system enforces convergent structure. Infinite scroll is not a new cognitive pathology; it is a crude, commercially motivated industrialization of an ancient learning strategy---one that captures the surface form of incremental distributed rehearsal while stripping the structural properties that make that strategy effective \cite{Goodhart1975}.

The access dimension deserves equal weight, and the critique systematically neglects it. The same device that delivers algorithmically optimized feeds also compresses what previously required physical travel, institutional affiliation, and sequential library access into a single portable interface, available in any context, in any language, at any moment. I have conducted research while traveling, accessed multilingual sources, and maintained scholarly work in conditions that would previously have required permanent institutional affiliation and geographic stability. This is not cognitive decline. It is a massive reduction in epistemic latency---capability expansion under compression.

\section{Epistemic Compression and the Collapse of Distance}

The critique of digital mediation systematically neglects a structural transformation that is at least as significant as any cost it introduces. Historically, access to distributed knowledge required physical travel, institutional affiliation, and sequential access to archives and expertise. These constraints imposed a sparse and path-dependent topology on knowledge acquisition. Contemporary digital systems collapse these distances, enabling simultaneous access to geographically and linguistically distributed sources without the friction of spatial separation.

This transformation has two inseparable consequences. It increases the size of the accessible option space, contributing to the stress dynamics analyzed in Section~7, and it enables forms of synthesis that were previously infeasible, as multiple domains can be brought into relation without the delays and costs imposed by spatial constraint \cite{Hutchins1995}. What appears, from within a scalar attentional model, as distraction is, from a structural perspective, the operation of a system whose epistemic latency has been dramatically reduced---a system capable of performing in parallel what previously required sequential institutional access spread across months and years.

\section{Event Stratification and the Recovery of Coherence}

The apparent opposition between fragmentation and coherence can be resolved by introducing an explicit stratification of the event space. In an unstructured feed, all inputs compete at the same level of immediacy. A notification, a conversational message, and a sustained argument occupy the same representational stratum and therefore produce interference, preventing the formation of stable trajectories \cite{Wiener1948}.

A stratified architecture assigns events to distinct layers with differentiated temporal and structural roles. Transient signals are separated from trajectory-defining commitments. Background monitoring is distinguished from foreground construction. History is preserved as an ordered, navigable sequence rather than an undifferentiated stream. Under such stratification, parallel streams are maintained without mutual collapse because they are not forced into the same representational layer. This provides a concrete mechanism for reconciling distributed attention with coherent long-range structure: not by reducing the number of active streams, but by organizing them into a structured temporal hierarchy in which each stream operates at the appropriate scale without interfering with the others.

\section{From Attention to Constraint: A Change of Primitive}

The final and deepest limitation of the continuous partial attention framework is that it treats attention as the primary explanatory primitive---as a quantity to be allocated, depleted, or captured \cite{Simon1969}. This choice of primitive obscures the underlying dynamics rather than illuminating them. What actually governs cognitive activity is not the distribution of attention per se but the resolution of constraints across representations \cite{Friston2010,Clark2013}.

Under a constraint-based view, attention is not a resource but an emergent property of constraint gradients. Cognitive systems move toward configurations that reduce inconsistency across active representations. What appears as distraction is often the system responding to unresolved constraints in adjacent contexts, where the gradient toward resolution is steeper than the gradient toward remaining in the current task. Notifications attract attention because they introduce new constraints or register unresolved states. Switching between tasks reflects the propagation of constraints across representational layers. Sustained focus corresponds to regions of the cognitive landscape where constraint satisfaction can proceed along a stable trajectory without being interrupted by higher-gradient unresolved structures elsewhere.

This reframing yields a unified interpretation of learning, work, and social interaction within a single framework. It also aligns with the broader direction in which cognitive science has moved toward understanding cognition as continuous constraint relaxation rather than discrete allocation of attentional quanta \cite{Friston2010}. Most importantly, it dissolves the puzzle of how multi-stream engagement can produce superior outcomes: when streams share structural relations, constraint resolution in one propagates to others, and the system converges faster than it would under sequential single-stream processing.

\section{The First Coordinate-System Error}

The argument developed in the preceding sections can now be stated with the precision it requires. Scryleus is observing real phenomena: distributed attention, increased cognitive load, altered social dynamics, subjective thinness of experience, elevated stress. These are genuine features of the current environment. The diagnostic failure is at the level of the model used to interpret them.

His framework operates on a scalar, single-threaded model of mind. Under that model, distributed attention appears as absence, multi-stream processing appears as fragmentation, altered social enforcement appears as decay, and subjective thinness appears as cognitive shallowing. The replacement framework---attention as a vector field, cognition as parallel constraint resolution, time as irreversible event stratification, stress as option-space expansion, thinness as projection artifact of dimensional compression---produces a different reading of every observation. Absence becomes multi-context embedding. Fragmentation becomes parallel constraint processing. Social change becomes enforcement reweighting under resource competition. Thinness becomes loss of representational degrees of freedom under dimensional reduction.

The observations are not disputed. The coordinate system for interpreting them is inadequate to the phenomena it purports to describe. What appears as absence is overextension across multiple coherent contexts, compressed into a medium that cannot represent it. The two interpretive systems are not merely in disagreement over conclusions; they preserve different invariants and are therefore incommensurable at the level of their fundamental representational commitments. Behaviors that are regular and productive in the field-stratified model must appear as anomalies in the scalar model, not because they are anomalous, but because the scalar model lacks the degrees of freedom required to encode them \cite{Barandes2023}.

\section{Metric Collapse and the Failure of Measurement}

The evaluation of technological systems relies on metrics that are assumed to faithfully represent underlying performance. Latency, engagement, conversion rates, and cost efficiency are treated as objective indicators of improvement \cite{Goodhart1975}. The central assumption is that these quantities preserve the structure of the phenomena they measure.

This assumption fails under dimensional reduction. When a high-dimensional process is projected into a scalar metric, information is not merely compressed but destroyed. Relationships between variables, temporal dependencies, and structural constraints are collapsed into a single value that cannot encode them. The metric becomes non-injective: multiple distinct system states map to the same measured outcome.

Formally, let $X$ denote the space of system states and $m \colon X \to \mathbb{R}$ a scalar metric. If $m$ is many-to-one, then there exist $x_1 \neq x_2$ such that $m(x_1) = m(x_2)$ despite $x_1$ and $x_2$ differing in properties that may be critical for system stability. Optimization over $m$ cannot distinguish between these states and will therefore select configurations that are optimal with respect to the metric while degrading unmeasured dimensions \cite{Ashby1956}.

This is not a secondary issue but a structural inevitability. Systems optimized under such metrics will systematically drift toward configurations that maximize measured performance while externalizing or suppressing the dimensions that the metric fails to capture. The resulting behavior appears as progress within the metric space and as degradation within the full system space. Goodhart's law captures the qualitative form of this phenomenon, but the deeper point is formal: it is not that metrics are gamed, but that non-injective projection makes gaming indistinguishable from genuine improvement \cite{Goodhart1975}.

The efficiency discourse surrounding artificial intelligence is therefore not merely incomplete but formally underdetermined. It evaluates systems using projections that cannot represent the invariants required for their stable operation. The resulting misalignment is not accidental; it is the direct and necessary consequence of optimizing within an insufficient coordinate system.

\section{The Externalization of Cost Under Efficiency Metrics}

The discourse surrounding artificial intelligence consistently evaluates systems through internal measures of performance: latency reduction, predictive accuracy, throughput, and cost efficiency. These metrics are treated as sufficient indicators of progress \cite{Goodhart1975}. What they fail to capture is that such efficiencies are not achieved in isolation but through the redistribution of cost across the system boundary.

The claim that improved prediction in domains such as insurance enables broader access at lower prices is a paradigmatic example. Increased predictive resolution does not produce equity; it produces segmentation. Individuals are no longer pooled into coarse risk categories but resolved into fine-grained profiles, allowing institutions to optimize inclusion and exclusion with increasing precision. The system does not become more just. It becomes more efficient at identifying which individuals are profitable to serve, and more precise in the determination of who falls beneath the threshold of inclusion.

The cost of this efficiency is not eliminated but externalized \cite{Prigogine1984}. It appears as pervasive data extraction, continuous surveillance, and the transformation of ordinary life into an input stream for predictive modeling \cite{Zuboff2019}. The individual does not simply purchase a service; they become a substrate upon which the system operates. The apparent reduction in price is offset by the loss of privacy, autonomy, and informational sovereignty. Efficiency, in this sense, is not a neutral optimization but a reallocation of thermodynamic and informational burden from institution to participant.

This pattern generalizes across the domains in which AI is being promoted as a democratizing force. In logistics, personalization, and automated service systems, efficiency gains correspond to an increased requirement that behavior be legible, predictable, and capturable. The system reduces its internal uncertainty by increasing the transparency of the user. What is framed as convenience is structurally identical to the reduction of degrees of freedom available to the individual within the system. The more efficient the system becomes in the internal sense, the narrower the space of viable deviation becomes for those whose behavior it models.

\section{Simulation and the Collapse of Irreversible History}

The transition from physical iteration to simulation-driven development is presented as a reversal of the traditional pipeline: learning precedes building. Thousands of candidate designs are evaluated in virtual environments, and only the optimal configuration is realized in material form. This is described as the elimination of waste. The deeper consequence is the elimination of irreversible history as a constitutive element of knowledge.

In a physical process, each attempt produces a trace that cannot be undone. Errors are not merely discarded; they become part of a trajectory that informs subsequent action \cite{Landauer1961}. This irreversibility grounds expertise, embedding knowledge in the geometry of lived interaction. Simulation replaces this structure with a space of reversible, non-binding possibilities. Paths explored within the simulation do not enter the shared historical record. They are neither persistent nor independently verifiable.

When such systems are coupled with ephemeral interfaces---feeds that cannot be replayed, histories that are selectively retained, and outputs that are generated on demand without stable provenance---the continuity of experience itself becomes fragmented. Information ceases to accumulate as a coherent trajectory and instead appears as a sequence of context-dependent renderings. The capacity to reconstruct how a conclusion was reached is lost, not because the process did not occur, but because it was never committed to a persistent, accessible record \cite{Wiener1948}.

From an entropic perspective, this introduces informational debt \cite{Shannon1948,Jaynes1957}. The system produces outputs whose generative histories are inaccessible, requiring trust in the system rather than verification through reconstruction. The user is positioned not as a participant in a shared unfolding of events, but as a consumer of precomputed results. The cost is the erosion of the very structure that allows knowledge to be grounded, shared, and contested.

\section{The Collapse of Skill Signaling Under Synthetic Saturation}

The claim that generative systems democratize creativity rests on the reduction of production costs. It assumes that lowering the barrier to creation increases access to participation. This assumption fails to account for the shift in the limiting factor from production to recognition.

When synthetic media can be generated at scale, the space of available artifacts expands dramatically. At the same time, the reliability of authenticity as a signal collapses. A visual, auditory, or textual artifact no longer encodes the history of its production in a way that can be trusted. A lifetime of skill development can be simulated within seconds, producing outputs indistinguishable from those generated through extended effort. The indistinguishability of authentic and synthetic artifacts is not a manageable risk to be addressed by governance mechanisms; it is a structural transformation in the epistemic function of evidence itself.

Under these conditions, the traditional pathway by which individuals without institutional support demonstrate competence---the public display of work---loses its discriminative power \cite{Zuboff2019}. The attention economy becomes saturated with artifacts detached from the histories that would otherwise validate them. Visibility becomes the primary scarce resource, and access to visibility is mediated by the very platforms that produce the saturation. For those who have historically depended on direct demonstration of skill as their only available form of credentialing, this is not an incidental disruption. It is the systematic elimination of the one remaining pathway that did not require prior institutional affiliation, advertising expenditure, or network access.

The consequence is not democratization but displacement. Those who previously relied on direct demonstration of skill are now competing within an environment where the signal of skill has been decoupled from its underlying process. The system does not elevate more participants; it reconfigures participation such that recognition depends on alignment with platform dynamics rather than the intrinsic properties of the work itself \cite{Goodhart1975}.

\section{The Surveillance Bottleneck and the Mediation of Reality}

The restructuring of informational pipelines has shifted the locus of control from distributed, persistent media to centralized, dynamically curated systems. Information no longer circulates as a stable artifact but as a mediated stream, subject to continuous filtering, ranking, and recomposition \cite{Zuboff2019}.

This produces what may be termed a surveillance bottleneck. In order to participate in the informational environment, all inputs must pass through systems that observe, record, and process them. The same infrastructure that enables access also enforces capture. There is no unmediated path between expression and reception. Every interaction is both communication and data acquisition. The claimed benefit of expanded access is inseparable from the structural requirement of total legibility.

As systems evolve toward real-time interpretation---devices that do not merely display information but explain it through overlays and contextual filtering---this mediation extends from communication into perception itself. The environment is no longer encountered directly but through a layer of interpretation that is both proprietary and adaptive. The distinction between representation and reality becomes increasingly difficult to maintain, as both are delivered through the same interface and both are subject to the same optimization criteria.

In such a system, the stability of history is contingent upon the policies and persistence guarantees of the mediating platform. Records can be modified, deprioritized, or removed \cite{Wiener1948}. Reconstruction of past states becomes dependent on access rights rather than inherent persistence. The shared field of reference required for collective verification is replaced by a set of overlapping but non-identical views, each shaped by the system's internal optimization criteria. A history that can be played back only through a commercial interface is not a history in any meaningful epistemic sense; it is an asset held in trust by a third party with no obligation to preserve it faithfully.

\section{Asymmetry of Visibility and Power}

The discourse of democratization assumes that increased access to tools produces a corresponding increase in agency. This assumption neglects a fundamental asymmetry: the difference between being visible to the system and being able to see the system.

Users operate under conditions of high observability. Their actions, preferences, and behavioral patterns are continuously recorded and analyzed. The system, by contrast, remains largely opaque. Its models, optimization criteria, and internal state transitions are not accessible to those subject to them \cite{Zuboff2019}. This produces a one-sided informational gradient that is not incidental but constitutive of the architecture's performance characteristics.

Let $V_u$ denote the visibility of the user to the system and $V_s$ the visibility of the system to the user. Contemporary platforms operate in a regime where
\[
V_u \gg V_s.
\]
This asymmetry enables precise modeling and manipulation of user behavior while preventing reciprocal understanding or intervention. The user participates in a system whose dynamics they cannot inspect, while the system continuously refines its model of the user. The appearance of empowerment---through access to tools and platforms---coexists with a structural reduction in the capacity to influence the conditions under which those tools operate.

This asymmetry is not incidental but constitutive of the current technological architecture. It is the condition under which predictive and optimization systems achieve their claimed performance. The informational gradient $V_u \gg V_s$ is not a design flaw to be corrected but a design feature to be maintained, because it is precisely this asymmetry that enables the extraction of behavioral data at the scale required for model training and targeted influence. Any account that treats increased access as equivalent to increased agency without accounting for this asymmetry is operating within an incomplete model of the system.

\section{Advertising as the Dominant Alignment Function}

The economic structure underlying contemporary digital systems is not incidental to their behavior; it is constitutive. The primary revenue model---the optimization of attention and behavior for commercial conversion---acts as the dominant alignment function for system design \cite{Goodhart1975,Zuboff2019}.

Within this framework, all system components are evaluated according to their contribution to engagement, retention, and influence. Personalization, recommendation, and generative capabilities are not neutral tools but mechanisms for increasing the efficiency with which user behavior can be shaped. The system does not merely respond to user preferences; it actively participates in their formation. What the dominant discourse frames as hyperpersonalization and expanded choice is, from within this structural analysis, the individualization and intensification of influence---the application of precision targeting to the formation of desire rather than merely its satisfaction.

This has two critical consequences. The first is that the optimization target is not truth, coherence, or user well-being, but measurable engagement. The second is that the system learns from and amplifies the very behaviors it induces, producing a feedback loop in which the environment becomes increasingly tuned to capture attention and reduce the viability of modes of interaction that do not align with conversion metrics. For individuals outside positions of economic privilege, this manifests as an environment saturated with targeted persuasion, where the space of available actions is continuously shaped by forces that operate at a level of precision and scale unavailable to the individual. The promise of personalization becomes the mechanism by which influence is individualized and intensified \cite{Zuboff2019}. Advertising is not an application of artificial intelligence; it is its primary organizing logic, and every other function of the system is developed in a context structured by that logic.

\section{The Geometry of Chokepoints}

The contemporary informational environment is not uniformly distributed but organized around a set of high-curvature regions through which the majority of communication and interaction must pass. These regions function as chokepoints: sites at which flow is concentrated, monitored, and regulated \cite{Wiener1948}.

In a geometric interpretation, the space of possible informational trajectories can be modeled as a manifold in which curvature corresponds to the concentration of control and influence. Regions of high curvature act as attractors, pulling trajectories toward them and reducing the diversity of viable paths. Formally, let $\mathcal{M}$ denote the manifold of informational states and $\kappa(x)$ its curvature at point $x$. Chokepoints correspond to regions where $\kappa(x)$ is large, producing geodesic convergence: independent trajectories that would otherwise remain distinct are forced into proximity, where they become subject to the same filtering, ranking, and optimization processes.

This concentration has two effects that are presented as separable but are structurally coupled. The first is efficient aggregation and processing of information, which is promoted as a benefit of scale. The second is a reduction in the system's resilience by making it dependent on a small number of structurally critical regions. Control over these regions confers disproportionate influence over the entire system. The surveillance bottleneck identified in the preceding section is the observational consequence of this geometric structure: all trajectories pass through monitored regions because the manifold has been shaped so that no low-curvature alternatives exist.

The critical structural problem is not the existence of chokepoints per se but the absence of alternative pathways that would allow information to circulate without passing through centralized control structures. A system in which all trajectories must pass through monitored regions cannot sustain independent verification or robust pluralism. It becomes a system in which participation is equivalent to movement through a controlled flow, and where the diversity of possible trajectories is systematically narrowed by the curvature of the underlying manifold.

\section{Informational Debt and Deferred Collapse}

Systems that suppress irreversibility and externalize cost do not eliminate these quantities; they defer their resolution. The resulting imbalance can be understood as informational debt: a discrepancy between the apparent coherence of system outputs and the unresolved structure required to support them \cite{Landauer1961,Shannon1948}.

Informational debt accumulates when outputs are generated without preserving the conditions of their production, when metrics are optimized at the expense of unmeasured variables, and when costs are displaced outside the system's accounting framework. Each such operation introduces latent inconsistency that must eventually be reconciled. Let $D(t)$ denote the accumulated informational debt at time $t$. In systems that continually externalize cost and suppress history, $D(t)$ is monotonically non-decreasing. The system remains operational so long as the debt can be deferred, but its stability decreases as $D(t)$ grows, because the gap between apparent and actual structural coherence widens.

Resolution of informational debt is not optional in the thermodynamic sense \cite{Prigogine1984}. It is not a catastrophic event but a phase transition: a reconfiguration in which previously suppressed constraints reassert themselves. This may take the form of loss of institutional trust, breakdown of shared epistemic reference, economic instability under conditions the model did not account for, or the failure of systems to generalize beyond the narrow conditions under which they were optimized. The critical point is that such transitions are not caused by external shock but by internal inconsistency that has been accumulated rather than resolved. Informational debt, like its thermodynamic counterpart, cannot be eliminated by accounting; it must be paid through the reintroduction of the very constraints that were previously suppressed.

\section{From Cognitive Misdiagnosis to Structural Misalignment}

The error identified in the continuous partial attention thesis---the use of a scalar model to interpret field-like cognition---reappears in the evaluation of technological systems themselves. Efficiency is treated as a scalar improvement, independent of the structural transformations required to produce it. Democratization is inferred from reduced cost, without accounting for the redistribution of control. Increased output is equated with increased participation, without examining the conditions under which that output is recognized and valued.

The result is a consistent pattern of misinterpretation. Systems that reorganize agency, compress history, and centralize control are described in terms of surface-level improvements. The deeper transformations are rendered invisible by the coordinate system used to evaluate them. Just as distributed attention appears as absence under a scalar model, distributed control appears as empowerment under a framework that lacks the degrees of freedom to represent mediation and dependency \cite{Ashby1956}.

A structurally adequate analysis must therefore operate at the level of constraints, trajectories, and fields. It must account for how systems shape the space of possible actions, how histories are recorded or erased, and how participation is mediated. Only within such a framework can the full set of consequences---both enabling and constraining---be made visible. The same framework that reveals multi-stream cognition as integration rather than fragmentation reveals AI-driven efficiency as redistribution rather than creation of value.

\section{Isomorphism Between Cognitive and Infrastructural Misdiagnosis}

The critique developed in this essay operates across two domains---cognition and technological infrastructure---and the claim that both exhibit the same coordinate-system error requires explicit formulation rather than repeated assertion.

Let $X_c$ denote the space of cognitive states and $X_s$ the space of system configurations. In both cases, a high-dimensional structure is projected into a scalar observable:
\[
m_c \colon X_c \to \mathbb{R}, \qquad m_s \colon X_s \to \mathbb{R}.
\]
In cognition, $m_c$ corresponds to scalar attention models that measure allocation to a single locus. In infrastructure, $m_s$ corresponds to efficiency metrics such as latency, engagement rate, or cost reduction. In both cases, the projection is non-injective. Distinct high-dimensional configurations map to identical scalar values, producing indistinguishability under the metric.

The consequences are formally parallel. In cognition, multi-layer constraint resolution appears as fragmentation because trajectories that are coherent at the level of $X_c$ are mapped to the same apparent attentional state under $m_c$. In infrastructure, redistribution of cost and centralization of control appear as efficiency because configurations that differ in their structural consequences for agency, history, and power are mapped to the same apparent performance value under $m_s$.

The isomorphism is therefore structural rather than causal. Both domains exhibit the composition
\[
X \xrightarrow{\;\mathcal{R}\;} X' \xrightarrow{\;m\;} \mathbb{R},
\]
where $\mathcal{R}$ is a dimensionality-reducing representation and $m$ a scalar metric. The error arises in both cases from treating $m$ as a faithful representation of $X$. The corrective move is identical in both cases: replace scalar evaluation with trajectory-and-field-based analysis. The implementations differ---in cognition, this entails interfaces that preserve multi-layer structure; in infrastructure, it entails systems that preserve history, expose cost, and resist collapse into single optimization targets---but the shared structure of the error implies a shared class of solutions defined by increased representational dimensionality.

\section{Structural Persistence of the Coordinate-System Error}

The persistence of scalar models and efficiency metrics is not solely an intellectual failure. It is structurally reinforced by the economic conditions under which contemporary systems operate, and understanding this reinforcement is necessary for understanding why the error is so stable.

Let $m$ denote a metric optimized by a system and let $\mathcal{D}$ denote the set of dimensions excluded from $m$. If optimizing $m$ produces gains that are immediately realizable, while degradation in $\mathcal{D}$ is delayed or externalized, then systems will converge toward configurations that maximize $m$ regardless of their impact on $\mathcal{D}$ \cite{Goodhart1975}. This dynamic is self-reinforcing. As systems become more dependent on $m$, alternative representations that capture $\mathcal{D}$ are marginalized. The coordinate system becomes not only a model but an infrastructure, embedded in measurement, evaluation, and decision-making processes that are themselves optimized to perpetuate it.

A shift to a higher-dimensional framework therefore requires more than conceptual revision. It requires the introduction of representations and metrics that internalize previously externalized dimensions. Without such changes, the system will continue to reproduce the same error, because the economic conditions that favor it remain in place. The persistence of the coordinate-system error is thus not an epistemic accident but a structural equilibrium: scalar models are maintained not because they are epistemically adequate but because they are economically stable. The coordinate-system error persists because it is profitable, not because it is correct.

\section{The Second Coordinate-System Error}

The critique can now be extended and unified. The first coordinate-system error misinterprets multi-layer cognition as fragmentation by projecting it through a scalar attention model that lacks the degrees of freedom to represent trajectory coherence. The second misinterprets system-level transformations as neutral improvements in efficiency and access by projecting them through a metric framework that lacks the degrees of freedom to represent externalized cost, historical fragility, and mediated agency.

In both cases, the error arises from the use of an insufficiently expressive model. Scalar metrics collapse multi-dimensional structure into single values, obscuring the redistribution of cost, the mediation of agency, and the degradation of historical continuity. Behaviors and systems that are structurally coherent within a higher-dimensional representation appear as anomalies or improvements when projected into a lower-dimensional one \cite{Barandes2023}. The two errors are not independent. They are instances of the same epistemological failure applied at different scales: the individual scale of cognition and the systemic scale of economic and informational infrastructure.

What appears as progress within the scalar frame may, within a structurally adequate frame, reveal itself as a reconfiguration of the conditions under which agency, knowledge, and participation are possible. The claims of the efficiency discourse are not false in isolation; they are expressed in a coordinate system that excludes the variables necessary to evaluate their full consequences. The omission of irreversibility, externalized cost, and field structure produces descriptions that are locally accurate and globally misleading.

\section{Pedagogical Reductionism and the Scalar Assessment Trap}

The coordinate-system error manifests with particular clarity in modern pedagogy, where standardized assessment functions as a scalar projection of a fundamentally trajectory-dependent process. Learning is treated as a state that can be sampled at a moment in time and assigned a numerical value, rather than as an irreversible transformation of a learner's internal representational structure.

Let $X$ denote the space of cognitive states and $\gamma \subset X$ a learning trajectory. Standardized testing defines a mapping $m_{\mathrm{edu}} \colon X \to \mathbb{R}$ which evaluates a learner at a single time $t$ without reference to $\gamma$. As this mapping is necessarily non-injective, distinct developmental trajectories that differ in depth, generality, and constraint integration may produce identical scores \cite{Simon1969}.

The result is a systematic preference for reversible knowledge. Strategies that maximize $m_{\mathrm{edu}}$ tend to minimize commitment: short-term memorization, pattern recognition without structural understanding, and procedural mimicry that leaves minimal trace in long-term cognitive organization. These strategies produce correct outputs while failing to encode the generative history required for transfer across domains. True mastery, by contrast, is characterized by irreversibility \cite{Landauer1961}. A concept, once integrated, constrains future reasoning across a wide class of problems. This corresponds to a transformation of the learner's constraint field $\mathcal{C}$, not merely an increase in stored information. The distinction is not quantitative but structural: mastery alters the geometry of the space in which reasoning occurs.

The reliance on scalar assessment therefore induces a form of educational informational debt. The system produces certified outputs---scores, credentials, rankings---without preserving the trajectory required to justify them. This debt becomes visible when learners are required to generalize beyond the narrow conditions under which the metric was optimized. A structurally adequate pedagogy would replace instantaneous evaluation with trajectory-sensitive measures that track the evolution of constraint resolution over time, the persistence of conceptual transformations, and the capacity to maintain coherence across increasing representational freedom. In such a framework, the kinds of parallel engagement documented in this essay---multi-domain reading, embodied learning alongside semantic intake, distributed incremental acquisition---would be interpreted not as fragmentation but as evidence of integration within a higher-dimensional cognitive field \cite{Hutchins1995}.

\section{Spatial Plenums and the Recovery of Interface Dimensionality}

The phenomenological thinness of contemporary digital experience is a direct consequence of dimensional reduction in interface design. High-dimensional cognitive processes are forced into sequential, flat representations that cannot preserve their relational structure \cite{Hutchins1995,Suchman2007}.

Let $\mathcal{I}$ denote the interface mapping from a cognitive state space $X$ to a presented representation. In current systems, $\mathcal{I}$ is effectively one-dimensional, organizing information as a time-ordered sequence. This imposes a total ordering on structures that are inherently partially ordered, eliminating spatial and relational degrees of freedom. Physical environments, by contrast, support volumetric indexing. Objects persist in space, allowing multiple contexts to coexist without interference. Partial states remain accessible, and relationships between elements can be encoded through proximity, orientation, and layering. These properties enable the maintenance of complex trajectories without repeated reconstruction. The transition to digital interfaces replaces this volumetric structure with streams: context is discarded when it leaves the viewport, and re-entry requires reconstruction from memory or search. This introduces cognitive cost not by excess complexity but by insufficient dimensionality.

A spatial plenum restores these degrees of freedom. Such a system treats information as embedded in a continuous field rather than a sequence, supporting persistent spatial indexing, layered stratification of temporal scales, and explicit visualization of trajectories. The user navigates not a feed but a manifold of states in which past configurations remain accessible and relational structure is preserved. The significance of this shift is not aesthetic but structural. By increasing the dimensionality of $\mathcal{I}$, the interface becomes capable of representing the trajectories that underlie cognition. This reduces informational debt by eliminating the need for repeated reconstruction and allows constraint propagation to occur across time and context \cite{Clark2013}. Distributed attention, in such an environment, becomes a stable configuration rather than a pathological deviation.

\section{Forensic Fragility and the Erosion of Verifiable History}

The stability of a shared epistemic environment depends on the persistence and accessibility of historical trajectories. In the absence of such persistence, outputs cannot be traced to their generative conditions, and verification becomes structurally impossible rather than merely difficult.

Contemporary digital systems increasingly operate in a regime where history is either compressed or inaccessible. Let $\mathcal{H}_t^*$ denote the full generative history of an artifact and $\mathcal{H}_t$ the history available to observers. In systems optimized for output rather than provenance, the mapping from $\mathcal{H}_t^*$ to $\mathcal{H}_t$ is lossy, producing $D(t) = \mathrm{dist}(\mathcal{H}_t, \mathcal{H}_t^*) > 0$. Synthetic media amplifies this condition catastrophically. When artifacts can be generated without encoding the trajectory of their production, the observable output ceases to function as evidence of its origin. The relationship between signal and source is severed, and the interpretive burden shifts from reconstruction to trust in the mediating platform.

This produces forensic fragility: a state in which truth is no longer grounded in the ability to replay and verify a sequence of events, but in the authority of systems that control access to representations of those events. The surveillance bottleneck intensifies this effect---all information flows through centralized systems that observe and record, yet do not provide symmetrical access to their records \cite{Zuboff2019}. Historical structure becomes partitioned and mediated, and participation in the informational environment requires acceptance of representations that cannot be independently verified.

The consequence for skill signaling identified earlier is now fully general. It is not merely that individual competence becomes unverifiable; it is that the epistemic function of evidence itself---the capacity to ground claims in reconstructible sequences of events---is degraded across all domains simultaneously. Restoring forensic integrity requires architectures that preserve irreversibility at the level of record: cryptographic provenance, append-only logs, and verifiable data lineage are not optional enhancements but structural requirements for maintaining a shared epistemic field \cite{Landauer1961}. Without such mechanisms, the system continues to accumulate informational debt, and the distinction between true and false collapses into a function of access rather than structure.

\section{Epistemic and Structural Criteria for Model Adequacy}

The argument advanced in this essay carries normative implications, but these arise from structural rather than moral considerations, and this distinction requires explicit statement to prevent the analysis from being read as a values dispute rather than an epistemological one.

A model is epistemically adequate if it preserves the distinctions required to reconstruct system trajectories. It is structurally adequate if optimization within the model does not produce degradation in dimensions that the system depends upon for stability. Scalar models fail both criteria. They collapse distinctions between trajectories, preventing reconstruction, and they permit optimization that degrades unmeasured dimensions. Field-stratified models, by contrast, preserve trajectory information and enable constraint propagation across time \cite{Friston2010,Wiener1948}.

The preference for higher-dimensional models is therefore not aesthetic or ethical in the first instance. It is grounded in their ability to represent and sustain the structures that systems require to remain coherent. Ethical and political implications follow from this, but they are not the primary justification. A system that cannot represent its own dynamics will produce systematically misleading conclusions about its own behavior. The correction is not to impose values on the system, but to adopt a representation in which its actual dynamics can be observed and understood. What this essay has identified as misdiagnosis---of cognition, of efficiency, of democratization---is not a failure of values but a failure of representational dimensionality. The political consequences are real and serious, but they are consequences of an epistemological failure, not its cause.

\section{Mechanisms of Irreversibility Suppression in AI Systems}

The loss of irreversible history in contemporary AI systems is not an abstract theoretical effect but the result of specific, identifiable architectural and economic choices that are individually rational under prevailing metrics and collectively destabilizing at the system level.

Interface design favors ephemerality. Content is presented in streams that are difficult to replay, search, or reconstruct. This eliminates persistent trajectories in favor of continuously refreshed states, and it does so not accidentally but because ephemeral interfaces reduce storage costs, increase re-engagement rates, and prevent users from accessing historical baselines against which current outputs could be compared. Data pipelines prioritize aggregation over provenance: inputs are collected, transformed, and incorporated into models without retaining accessible mappings between outputs and their originating data. The result is a function that is evaluable but not invertible in practice \cite{Shannon1948}.

Generative systems are optimized for output quality rather than traceability. A generated artifact does not encode the sequence of transformations that produced it, nor does it provide access to the latent states traversed during generation. This is not a technical limitation but an optimization choice: traceability increases computational cost and creates legal and reputational exposure. Finally, economic incentives reinforce all of these design choices simultaneously. Systems that preserve detailed history, expose internal structure, or enable reconstruction are more expensive to operate, harder to defend against liability claims, and less amenable to the kind of behavioral optimization that drives advertising revenue. These mechanisms collectively produce the suppression of irreversibility identified in the theoretical analysis. They are not incidental features but structurally aligned with the optimization targets of the system under current economic conditions \cite{Goodhart1975,Zuboff2019}.

\section{Irreversibility, Entropy, and the Reconstruction of the Plenum}

The argument developed across this essay can now be expressed in a unified formal language. What has appeared, at different points, as a critique of attention theory, digital interface design, and artificial intelligence discourse is in fact a single structural diagnosis: the systematic misrepresentation of irreversible processes within frameworks that treat systems as reversible, memoryless, and scalar \cite{Prigogine1984,Landauer1961}.

The central invariant that has been lost across these domains is irreversibility. In cognition, it appears as the failure to recognize that coherence is a property of trajectories rather than instantaneous states. In technological systems, it appears as the replacement of persistent history with ephemeral feeds and simulation outputs that do not retain the path by which they were generated. In economic structures, it appears as the externalization of cost into domains that are not recorded within the system's internal accounting.

Let $\mathcal{H}$ denote the space of histories and $\Omega_t$ the space of available possibilities at time $t$. A system that preserves irreversibility evolves according to a monotonic expansion of $\mathcal{H}$ and a corresponding contraction or restructuring of $\Omega_t$. Each committed action reduces the space of future possibilities while increasing the structure of the past. Coherence arises from the compatibility of these accumulated commitments \cite{Wiener1948,Ashby1956}.

By contrast, a system that suppresses irreversibility---through simulation, ephemeral interfaces, or non-persistent mediation---decouples $\mathcal{H}$ from observable outputs. Actions appear without reconstructible provenance. Possibilities are explored without commitment. The system presents the illusion of a continuously refreshed present, in which outcomes are delivered without visible causal history. Formally, this corresponds to a projection
\[
(\Omega_t,\, \mathcal{H}_t) \;\longrightarrow\; \Omega_t',
\]
in which the history component is either discarded or rendered inaccessible. The result is a loss of constraint propagation across time: the system no longer accumulates structure in a way that can be inspected, verified, or meaningfully extended \cite{Landauer1961}.

This loss can be understood in entropic terms \cite{Jaynes1957,Shannon1948}. Let $S(\mathcal{H}_t)$ denote the entropy associated with the system's historical record, interpreted as the degree of unresolved contradiction or uncertainty embedded within it. In a well-formed system, entropy is reduced locally through the accumulation of coherent structure, even as the global system evolves. Constraint resolution produces ordered trajectories that can be revisited and extended. In the absence of persistent history, entropy is not resolved but displaced. It reappears as uncertainty about provenance, as instability in interpretation, and as dependence on external authority to validate outputs. The system becomes informationally lossy: it produces results that cannot be traced back to their generating conditions. This is the structural origin of the informational debt identified in the analysis of simulation-driven development. The system has performed computation without retaining the structure necessary to justify its conclusions.

The concept of a plenum provides the appropriate corrective abstraction. A plenum is not a collection of isolated states but a continuous field in which scalar quantities, vector flows, and entropic distributions interact to produce observable structure \cite{Feynman1963,Barandes2023}. Within such a field, no event is independent of its context, and no outcome is detached from its generating history. The state of the system at any point is a function of its entire trajectory, not merely its current configuration. This is the formal counterpart to the observation that coherence is longitudinal rather than instantaneous.

Recasting the systems under consideration within this framework yields a different set of design principles. Interfaces must preserve trajectory: information should not be presented as an unstructured stream but as a navigable history in which causal relations are maintained. Computation must be accountable: outputs should carry with them the structure of their generation, allowing reconstruction and verification. Optimization must be constrained by global coherence: local efficiency gains that increase systemic entropy or externalize cost must be recognized as destabilizing rather than progressive \cite{Prigogine1984}.

These principles are not normative in the moral sense; they are structural requirements for the maintenance of coherence in systems that operate over time. A system that erases its own history cannot sustain stable knowledge. A system that externalizes its costs cannot maintain equilibrium. A system that collapses multi-dimensional structure into scalar metrics will misinterpret its own behavior.

The critique of contemporary technological discourse can therefore be restated with precision. It is not that the claims of increased efficiency, accessibility, or capability are false in isolation. It is that they are expressed in a coordinate system that excludes the variables necessary to evaluate their consequences. The omission of irreversibility, entropy, and field structure produces descriptions that are locally accurate but globally misleading.

Restoring these variables reveals a different picture. The expansion of computational capability coincides with the compression of historical continuity. The reduction of production cost coincides with the saturation of recognition channels. The increase in informational access coincides with the centralization of mediation. These are not contradictions but coupled transformations within a single field. The forward-looking claim that interfaces must evolve toward spatial, embodied, and stratified architectures is not a preference but a structural necessity: only within such architectures can the trajectory structure of cognition, the persistence of historical record, and the accountability of computation coexist.

The task, then, is not to reject the systems that produce these effects but to reconfigure them so that their operation remains embedded within a coherent plenum---a shared, verifiable, non-ephemeral field of historical structure. Without such embedding, the system will continue to generate outputs that are locally optimized and globally incoherent, accumulating informational debt that must eventually be resolved, and misrepresenting as progress the progressive erosion of the conditions under which knowledge, agency, and participation remain possible.

\section{Toward a High-Dimensional Framework of Agency and History}

The failures examined across cognition, pedagogy, interface design, epistemic infrastructure, and economic organization are instances of a single structural error: the reduction of trajectory-dependent systems to scalar representations. In each domain, a high-dimensional process is projected into a coordinate system that cannot represent its invariants, and the resulting distortions---fragmentation, efficiency illusions, skill-signal collapse, forensic instability, and pedagogical debt---are not independent phenomena but consequences of optimization within an insufficient representation.

The corrective is not the rejection of technology or of measurement, but the adoption of coordinate systems capable of preserving structure. This entails increasing representational dimensionality, maintaining access to historical trajectories, and designing systems in which constraint propagation remains possible across time and context. The concept of the plenum provides a unifying criterion for this project. Systems that sustain a plenum preserve history, expose structure, and allow participants to navigate and reconstruct trajectories. Systems that suppress it replace structure with output, trajectory with state, and verification with mediated access.

The choice between these regimes is not merely technical. It determines whether human activity remains embedded in a shared, navigable field of history or is reduced to interaction with precomputed outputs whose origins cannot be recovered. A high-dimensional framework does not eliminate complexity; it renders it visible and tractable. In doing so, it restores the possibility of coherence in systems that would otherwise converge toward distortion under the pressure of scalar optimization.

The central theorem of the following appendix formalizes this conclusion. Metric optimization under non-injective projection is not merely likely to produce distortion; it is provably required to accumulate informational debt in expectation. The misdiagnosis of distributed cognition and the misinterpretation of technological efficiency are therefore not contingent failures that better intentions could correct. They are structural outcomes of operating within coordinate systems that exclude the variables necessary to represent what is actually taking place. Changing the outcome requires changing the coordinate system.

\appendix

\section{Formal Primitives of the Field-Stratified Framework}

This appendix introduces a minimal set of formal primitives sufficient to express the structural claims developed throughout the essay. These definitions are not intended as a complete theory but as a coordinate system in which the phenomena under discussion can be represented without loss of essential structure.

\subsection*{Definition 1 (Plenum)}

A \emph{plenum} is a continuous field $\mathcal{P}$ over a domain $D \subseteq \mathbb{R}^n$ equipped with three coupled components:
\[
\mathcal{P}(x,t) = \bigl(\Phi(x,t),\; \mathbf{v}(x,t),\; S(x,t)\bigr),
\]
where $\Phi$ is a scalar potential, $\mathbf{v}$ is a vector flow field, and $S$ is an entropy density. The evolution of $\mathcal{P}$ is governed by constraint propagation across these components, such that local changes in any component induce changes in the others. The state of the system at $(x,t)$ is therefore a function of its trajectory through the field, not merely its instantaneous configuration.

\subsection*{Definition 2 (Irreversible Trajectory)}

Let $\mathcal{H}_t$ denote the history of a system up to time $t$. An \emph{irreversible trajectory} is a sequence of states $\gamma = (x_0, x_1, \ldots, x_t)$ such that there exists no inverse mapping that reconstructs $\gamma$ from its terminal state $x_t$ alone. Irreversibility implies that information about the sequence of transitions is not recoverable from the final state without access to $\mathcal{H}_t$. Coherence is defined over $\gamma$, not over $x_t$.

\subsection*{Definition 3 (Constraint Field)}

A \emph{constraint field} $\mathcal{C}$ over a state space $X$ assigns to each state $x \in X$ a scalar value representing unresolved inconsistency:
\[
\mathcal{C} \colon X \to \mathbb{R}_{\geq 0}.
\]
System evolution corresponds to movement along trajectories that locally reduce $\mathcal{C}$. Attention, in this framework, is not a resource but an emergent property of gradients in $\mathcal{C}$:
\[
\mathbf{a}(x) \sim -\nabla \mathcal{C}(x).
\]

\subsection*{Definition 4 (Metric Projection and Collapse)}

Let $X$ be a high-dimensional state space and $m \colon X \to \mathbb{R}$ a scalar metric. A \emph{metric projection} is the mapping induced by $m$. A projection exhibits \emph{metric collapse} if it is non-injective:
\[
\exists\, x_1 \neq x_2 \in X \;\text{ such that }\; m(x_1) = m(x_2).
\]
Under metric collapse, optimization over $m$ cannot distinguish between structurally distinct states, leading to systematic misalignment between measured performance and actual system behavior.

\subsection*{Definition 5 (Chokepoint Curvature)}

Let $\mathcal{M}$ be a manifold of informational states with curvature function $\kappa \colon \mathcal{M} \to \mathbb{R}_{\geq 0}$. A region $U \subset \mathcal{M}$ is a \emph{chokepoint} if $\kappa(x)$ is locally maximized over $U$, producing geodesic convergence:
\[
\forall\, \gamma_1, \gamma_2 \in \Gamma,\quad \gamma_1 \neq \gamma_2 \;\Rightarrow\; d\bigl(\gamma_1(t), \gamma_2(t)\bigr) \to 0 \;\text{ as }\; t \to t_U.
\]
Chokepoints concentrate flow, enabling control and observation at the cost of reduced trajectory diversity.

\subsection*{Definition 6 (Informational Debt)}

Let $\mathcal{H}_t$ denote the accessible history of a system and $\mathcal{H}_t^*$ the full generative history required to reconstruct its outputs. The \emph{informational debt} at time $t$ is:
\[
D(t) = \mathrm{dist}(\mathcal{H}_t,\, \mathcal{H}_t^*),
\]
where $\mathrm{dist}$ measures the discrepancy between accessible and required historical structure. A system accumulates informational debt when it produces outputs without preserving sufficient history to reconstruct their derivation. Increasing $D(t)$ corresponds to decreasing epistemic stability.

\subsection*{Definition 7 (Externalized Cost)}

Let $C_{\mathrm{int}}$ denote the cost measured within a system and $C_{\mathrm{ext}}$ the cost imposed outside its accounting boundary. A system exhibits \emph{externalization} if:
\[
C_{\mathrm{total}} = C_{\mathrm{int}} + C_{\mathrm{ext}}, \quad \text{with} \quad C_{\mathrm{ext}} \notin \operatorname{argmin}(C_{\mathrm{int}}).
\]
Optimization over $C_{\mathrm{int}}$ alone produces configurations that reduce internal cost while increasing total cost. Efficiency, under such conditions, is a projection artifact rather than a global improvement.

\subsection*{Operationalization}

These quantities admit empirical approximation and are not merely conceptual. The effective degrees of freedom of a representation are given by $\mathrm{DoF}_{\mathrm{eff}} = \mathrm{rank}(\mathcal{R})$ where $\mathcal{R}$ is the representational mapping; loss of representational capacity corresponds to a reduction in this rank. Constraint fields can be approximated through inconsistency measures across representations: given a set of representations $\{r_i\}$, one may define $\mathcal{C} = \sum_{i,j} d(r_i, r_j)$ where $d$ measures incompatibility. Informational debt can be operationalized through reconstruction error: letting $f$ denote a generative process and $\hat{f}$ a reconstruction derived from accessible history, $D = \mathbb{E}[\|f^{-1}(y) - \hat{f}^{-1}(y)\|]$ where $y$ is an observed output. These correspond to measurable properties of real systems: rank deficiency in learned representations, divergence across predictive models, reconstruction error in generative pipelines, and loss of provenance in data systems.

\section{Theorem: Metric Optimization Under Dimensional Reduction Produces Informational Debt}

The shared structure of the cognitive and infrastructural failures identified in the main text admits formal expression. We state the central result of the paper.

\begin{theorem}[Dimensional Reduction and Informational Debt]
Let $X$ be a high-dimensional state space with trajectories $\gamma \subset X$, and let $m \colon X \to \mathbb{R}$ be a non-injective scalar metric. Consider a system that evolves according to
\[
x_{t+1} = \operatorname{argmin}_{x \in X} m(x)
\]
at each step, without preserving full trajectory history $\mathcal{H}_t$. Then, under repeated optimization of $m$, the system will \textup{(i)} converge toward states that are optimal with respect to $m$ but not necessarily with respect to the full structure of $X$, and \textup{(ii)} accumulate informational debt $D(t) = \mathrm{dist}(\mathcal{H}_t, \mathcal{H}_t^*)$ that is monotonically non-decreasing in expectation.
\end{theorem}

\subsection*{Proof Sketch}

The proof proceeds in two steps. First, because $m$ is non-injective, there exist distinct states $x_1, x_2 \in X$ such that $m(x_1) = m(x_2)$ while $x_1$ and $x_2$ differ in dimensions not represented by $m$. Optimization over $m$ cannot distinguish between these states; the selection process therefore admits trajectories that preserve optimality under $m$ while degrading structure in the null space of $m$.

Second, the absence of preserved history implies that the system cannot reconstruct which trajectory produced a given state. Each step that selects among indistinguishable states without recording the distinguishing dimensions increases the discrepancy between $\mathcal{H}_t$ and $\mathcal{H}_t^*$. For each step $t$, there exists a non-zero probability that the system transitions to a state whose distinguishing information lies outside the range of $m$ and is not recorded in $\mathcal{H}_t$. This produces an expected increase in reconstruction error:
\[
\mathbb{E}[D(t+1)] \geq \mathbb{E}[D(t)].
\]
Thus, optimization under a non-injective metric without trajectory preservation necessarily produces accumulation of informational debt. \hfill $\square$

\begin{corollary}[Cognitive Misdiagnosis]
Let $X_c$ denote the space of cognitive trajectories and $m_c$ a scalar attention metric. Then multi-layer constraint resolution processes that are distinct in $X_c$ may be mapped to identical values under $m_c$, producing apparent fragmentation despite underlying coherence.
\end{corollary}

\begin{corollary}[Efficiency Misinterpretation]
Let $X_s$ denote the space of system configurations and $m_s$ an efficiency metric. Then configurations that differ in externalized cost, historical persistence, or agency distribution may be indistinguishable under $m_s$, producing apparent efficiency gains that correspond to redistribution of unmeasured cost.
\end{corollary}

\begin{corollary}[Design Constraint]
Any system that optimizes a non-injective scalar metric without preserving trajectory information will exhibit increasing divergence between measured performance and structural coherence. Avoiding this requires either \textup{(i)} increasing the dimensionality of the metric or \textup{(ii)} preserving sufficient history to reconstruct trajectories.
\end{corollary}

\noindent These three corollaries formalize the central claim of the essay. The misdiagnosis of distributed cognition and the misinterpretation of technological efficiency are not independent errors. They are instances of the same structural phenomenon: optimization within a coordinate system that cannot represent the invariants required for coherent system behavior. The theorem establishes that this outcome is not contingent but provably necessary given the conditions. The solution---in cognition, in interface design, in infrastructure, in economic accounting---is the same in every case: increase the dimensionality of the representation until it is adequate to the system being described.

\bibliographystyle{plain}
\begin{thebibliography}{99}

\bibitem{Shannon1948}
C.~E. Shannon,
``A Mathematical Theory of Communication,''
\emph{Bell System Technical Journal}, vol.~27, pp.~379--423, 1948.

\bibitem{Wiener1948}
N.~Wiener,
\emph{Cybernetics: Or Control and Communication in the Animal and the Machine},
MIT Press, 1948.

\bibitem{Ashby1956}
W.~R. Ashby,
\emph{An Introduction to Cybernetics},
Chapman \& Hall, 1956.

\bibitem{Simon1969}
H.~A. Simon,
\emph{The Sciences of the Artificial},
MIT Press, 1969.

\bibitem{Feynman1963}
R.~P. Feynman,
\emph{The Feynman Lectures on Physics},
Addison-Wesley, 1963.

\bibitem{Prigogine1984}
I.~Prigogine and I.~Stengers,
\emph{Order Out of Chaos: Man's New Dialogue with Nature},
Bantam Books, 1984.

\bibitem{Landauer1961}
R.~Landauer,
``Irreversibility and Heat Generation in the Computing Process,''
\emph{IBM Journal of Research and Development}, vol.~5, no.~3, pp.~183--191, 1961.

\bibitem{Jaynes1957}
E.~T. Jaynes,
``Information Theory and Statistical Mechanics,''
\emph{Physical Review}, vol.~106, pp.~620--630, 1957.

\bibitem{Friston2010}
K.~Friston,
``The Free-Energy Principle: A Unified Brain Theory?''
\emph{Nature Reviews Neuroscience}, vol.~11, pp.~127--138, 2010.

\bibitem{Clark2013}
A.~Clark,
\emph{Surfing Uncertainty: Prediction, Action, and the Embodied Mind},
Oxford University Press, 2013.

\bibitem{Hutchins1995}
E.~Hutchins,
\emph{Cognition in the Wild},
MIT Press, 1995.

\bibitem{Suchman2007}
L.~Suchman,
\emph{Human-Machine Reconfigurations},
Cambridge University Press, 2007.

\bibitem{Zuboff2019}
S.~Zuboff,
\emph{The Age of Surveillance Capitalism},
PublicAffairs, 2019.

\bibitem{Goodhart1975}
C.~A.~E. Goodhart,
``Problems of Monetary Management: The U.K. Experience,''
in \emph{Papers in Monetary Economics}, 1975.

\bibitem{Barandes2023}
J.~Barandes,
\emph{The Science of What Is Possible: A New Look at Reality Through Quantum Theory},
Oxford University Press, 2023.

\end{thebibliography}

\end{document}

